% !TEX program = xelatex
% A quick reference built with the templates themselves: source on the left,
% what it prints on the right. Every snippet is compiled right here.
\documentclass[lang=zh, thmnames=native, letters=bold, tikz=full, twocolumn]{akizumi-note}
\usepackage{akizumi-logic}
\geometry{landscape, margin=1.1cm, top=1.55cm, bottom=1.25cm, headsep=0.45cm}
\setlength{\columnsep}{0.75cm}
\raggedbottom
\akizumisetup{course=akizumi-latex 速查表, headright={左：原始碼\quad 右：結果}}
\title{akizumi-latex 速查表}

% ---------- look of the sheet ----------
\titleformat{\section}[block]{\sffamily\normalsize\bfseries\color{accent}}{\thesection}{0.6em}{}
  [\nobreak\vspace{-0.6ex}{\color{rule}\nobreak\titlerule[0.6pt]}\nobreak]
\titlespacing*{\section}{0pt}{0.9em}{0.45em}
\setlist{itemsep=0pt, topsep=0.2em}
\renewcommand{\arraystretch}{1.04}
\AtBeginDocument{\setlength{\headwidth}{\textwidth}}  % full width after landscape
\newcommand{\csfont}{\ttfamily\fontsize{7.6}{9.4}\selectfont}
\newlength{\cscolumn}
\AtBeginDocument{\settowidth{\cscolumn}{\csfont 0}}

% Source in the code font: control sequences in accent, comments muted.
% Fixed columns: xeCJK sets a CJK character in two of them, so comments
% aligned in this file (CJK counted twice) line up in the PDF too. $ and _
% are not letters, so \abs in "$\abs" is a control sequence like any other.
\lstdefinestyle{cs}{style=akizumi, language={}, numbers=none,
  basicstyle=\csfont, columns=fixed, basewidth=\cscolumn,
  alsoother={\$_}, alsoletter={\\}, morekeywords={\\begin}, keywordsprefix=\\,
  keywordstyle=\color{accent}, morecomment=[l]\%, commentstyle=\color{mute}, breakindent=1em}
\tcbset{cs/base/.style={enhanced, breakable=false, colframe=rule, boxrule=0.5pt, arc=1.5pt,
    left=5pt, right=5pt, top=4pt, bottom=4pt,
    before skip=0.45em, after skip=0.45em,
    listing options={style=cs}, listing file={\jobname-cs.tex}}}
% \begin{cs}[ratio] ... \end{cs}: source left, result right
\DeclareTCBListing{cs}{ !O{0.52} }{cs/base, listing side text, colback=white,
  lefthand ratio=#1, sidebyside gap=10pt, segmentation hidden,
  sidebyside align=top seam, fontlower=\small\raggedright,
  underlay={\begin{tcbclipinterior}\fill[wash] (interior.south west) rectangle
    ([xshift=-5pt]segmentation.north);\end{tcbclipinterior}}}
% \begin{cssrc}{note} ... \end{cssrc}: source only (other document kinds)
\DeclareTCBListing{cssrc}{ m }{cs/base, listing only, colback=wash,
  overlay={\node[anchor=north east, inner sep=3pt, font=\sffamily\scriptsize,
    text=mute] at (frame.north east) {#1};}}

% \mrow{source}: the source in code font, then what it prints
\ExplSyntaxOn
\tl_new:N \l__cs_tl
\NewDocumentCommand \mrow { v }
  {
    {\csfont #1} &
    \tl_set_rescan:Nnn \l__cs_tl { } {#1} $\l__cs_tl$
  }
\ExplSyntaxOff
\newcommand{\opt}[1]{{\csfont #1}}
% the default value; \leavevmode: a \color first in a p cell adds a \parskip
\newcommand{\dflt}[1]{\leavevmode{\csfont\color{accent}#1}}
% \csnote{...}: a muted line under a box or table
\newcommand{\csnote}[1]{\par{\footnotesize\color{mute}#1\par}}
% A heading right after a line of text (a \csnote; the last node is the line
% or the colour's whatsit, not the glue a box leaves) would get the \parskip
% on top of its own space, more than after a box: take it back there.
\NewCommandCopy{\cssection}{\section}
\RenewDocumentCommand{\section}{}{\par
  \ifnum\lastnodetype>0 \ifnum\lastnodetype<11 \vspace{-\parskip}\fi\fi\cssection}

\begin{document}
\twocolumn[{%
  \sffamily
  {\LARGE\bfseries akizumi-latex 速查表}\hfill
  {\small\color{mute}每段左邊是原始碼，右邊是印出來的樣子；完整說明：\texttt{ax docs 主題} 或 \texttt{docs/reference.md}}\par
  \vspace{0.4em}{\color{rule}\rule{\linewidth}{0.6pt}}\par\vspace{0.6em}}]

\footnotesize
\section{文件種類}\nobreak
\begin{tabularx}{\linewidth}{@{}l l Y@{}}
\toprule
\opt{ax new} & \opt{\textbackslash documentclass} & 用途 \\
\midrule
\opt{hw}     & \opt{akizumi-hw}   & 作業：\opt{problem}、\opt{parts}、\opt{solution} \\
\opt{note}   & \opt{akizumi-note} & 講義、筆記、報告（這份就是） \\
\opt{exam}   & \opt{akizumi-exam} & 考卷：exam 的 \opt{\textbackslash question[10]}、\opt{\textbackslash gradetable} \\
\opt{slides} & \opt{beamer} + \opt{\textbackslash usetheme\{Akizumi\}} & 簡報 \\
\opt{book}   & \opt{akizumi-book} & 論文、長篇講義：章、部、封面 \\
\opt{paper}  & 會議模板 + \opt{akizumi-math} & 投稿論文（pdfLaTeX） \\
\bottomrule
\end{tabularx}


\section{ax 指令}\nobreak
\begin{tabularx}{\linewidth}{@{}>{\raggedright\arraybackslash}p{13.5em} Y@{}}
\toprule
指令 & 作用 \\
\midrule
\opt{ax new hw hw3 --course 微積分} & 建立專案；不給參數會用選單問 \\
\opt{ax} & 專案狀態：種類、模板版本、PDF 更新時間 \\
\opt{ax build}\quad\opt{--chapter 3} & 編譯一次；失敗時指出哪一行 \\
\opt{ax watch}\quad\opt{--bg --log --stop} & 存檔就重編，可以在背景跑；\opt{--status} 看結果 \\
\opt{ax both}\quad\opt{ax open --solutions} & 題目卷和解答各一份；打開解答（\opt{ax open} 打開 PDF） \\
\opt{ax count}\quad\opt{--limit 3000} & 字數：英文字加中日韓字 \\
\opt{ax pack}\quad\opt{ax freeze} & 打包上傳 Overleaf；把模板複製進專案 \\
\opt{ax upgrade}\quad\opt{ax --version} & 換到最新的模板版本；最新的版本是哪個 \\
\opt{ax migrate old.tex} & 舊文件轉成 akizumi 版的副本 \\
\opt{ax docs 定理}\quad\opt{--cheatsheet} & 在終端機查參考手冊；打開這份速查表 \\
\opt{ax doctor}\quad\opt{ax config} & 出問題先跑；設定檔在哪裡 \\
\opt{ax completion zsh} & Tab 補全腳本（也有 bash、powershell） \\
\opt{ax selftest} & 改了模板之後：重編範例、測試、手冊片段 \\
\midrule
\opt{.ax.json} & \opt{"prebuild"}：編譯前先跑的指令；\opt{"shell\_escape": true}：LaTeX 可以執行外部程式；\opt{"template": "latest"}：用模板目前的內容（手冊裡比 v0.8 新的功能） \\
\bottomrule
\end{tabularx}

\section{常用的 class 選項}\nobreak
\begin{tabularx}{\linewidth}{@{}>{\raggedright\arraybackslash}p{6.4em} >{\raggedright\arraybackslash}p{12.2em} Y@{}}
\toprule
選項 & 值（\textcolor{accent}{藍色}是預設） & 作用 \\
\midrule
\opt{lang}       & \dflt{en} \opt{zh ja ko} & 主語言：日期、圖表名稱、字型 \\
\opt{thmnames}   & \dflt{english} \opt{native} & Theorem／定理、Proof／證明 \\
\opt{text}       & \dflt{serif} \opt{sans times palatino} & 內文字型；標題一律 Lato \\
\opt{paper}      & \dflt{a4} \opt{letter pad} & \opt{pad}：3:4 平板版面 \\
\opt{layout}     & \dflt{normal} \opt{margin} & \opt{margin}：寬邊欄放 \opt{\textbackslash sidenote} \\
\opt{byline}     & \dflt{false} \opt{true} & 標題下印「作者 · 學號 · 日期」 \\
\opt{fontsize}   & \dflt{11pt} \opt{12pt 13pt 14pt} & 字級；也可以只寫 \opt{13pt} \\
\opt{par}        & \dflt{skip} \opt{indent} & \opt{indent}：首行縮排 2 字 \\
\opt{secnum}     & \dflt{arabic} \opt{zh} & \opt{zh}：一、（一）1. \\
\opt{kai}        & （寫上就開） & 中文內文用楷書 \\
\opt{bib}        & 檔名，例如 \opt{refs.bib} & 參考文獻（biber） \\
\opt{bibstyle}   & \dflt{numeric} \opt{authoryear apa} & 引用格式 \\
\opt{cite}       & \dflt{normal} \opt{super} & \opt{super}：上標數字 \\
\opt{index}      & \dflt{false} \opt{true} & 索引、符號列表 \\
\bottomrule
\end{tabularx}

\begin{cssrc}{note}
\documentclass[lang=zh, thmnames=native, layout=margin]{akizumi-note}
\akizumisetup{course=常微分方程, kicker={ODE · 第 3 講}}
\title{一階線性方程}  \author{Aaron Lin}
\begin{document}
\maketitle
\end{cssrc}

\section{定理、程式碼、數學、圖的選項}\nobreak
\begin{tabularx}{\linewidth}{@{}>{\raggedright\arraybackslash}p{6.4em} >{\raggedright\arraybackslash}p{12.2em} Y@{}}
\toprule
選項 & 值（\textcolor{accent}{藍色}是預設） & 作用 \\
\midrule
\opt{thmcounter} & \dflt{shared} \opt{grouped split} & 定理類共用編號、分三類、各自編號 \\
\opt{thmwithin}  & \opt{section chapter none} & 編號前加節號或章號；預設講義 \opt{section}、書 \opt{chapter}、其他 \opt{none} \\
\opt{asides}     & \dflt{numbered} \opt{unnumbered} & 例、註不編號 \\
\opt{aliases}    & \dflt{none} \opt{short} & 短名稱 \opt{thm} \opt{lem} \opt{defn}… \\
\opt{thmname}    & \opt{\{corollary=系理\}} & 換掉某一種的名稱 \\
\opt{solutions}  & \opt{inline chapter end hide}；\opt{true false} & 解答放哪裡；預設講義和書 \opt{inline}、作業 \opt{true}、考卷 \opt{false} \\
\opt{exercisewithin} & \opt{chapter section none} & 習題的編號單位，預設同 \opt{thmwithin} \\
\opt{code}       & \dflt{listings} \opt{minted} & 程式碼上色引擎 \\
\opt{verbatim}   & \dflt{plain} \opt{styled} & \opt{verbatim} 環境也用程式碼的灰底 \\
\opt{letters}    & \opt{bold cal bb greek all single} & 字母捷徑 \opt{\textbackslash va}、\opt{\textbackslash sA}…；預設不開 \\
\opt{delims}     & \dflt{manual} \opt{auto} & \opt{auto}：括號不加星號就自動變大 \\
\opt{bb}         & \dflt{light} \opt{regular ams} & 黑板粗體的粗細 \\
\opt{units}      & \dflt{false} \opt{true} & siunitx：\opt{\textbackslash qty\{9.8\}\{\textbackslash metre\}} \\
\opt{tikz}       & \dflt{basic} \opt{full} & 更多 TikZ 函式庫 \\
\opt{plots}      & \dflt{false} \opt{true} & pgfplots，配色跟模板一致 \\
\bottomrule
\end{tabularx}

\section{定理、證明、交叉引用}
\begin{cs}
\begin{definition}[Collapse]\label{def:c}
若 $\Var[f(\vb{o})] = 0$，稱 $f$ collapse。
\end{definition}
\begin{theorem}\label{thm:a}
$\abs{x} \le 1 \implies x^2 \le 1$。
\end{theorem}
\begin{proof}
兩邊平方。
\end{proof}
見\cref{thm:a}、\cref{def:c}。
\end{cs}

\begin{cs}
\begin{restatable}[Cram\'er--Wold]
  {theorem}{cw}\label{thm:cw}
$\vb{z} \sim \Normal{\vb{0}}{\vb{I}}$ 若且唯若
每個 $\vb{u}\T\vb{z}$ 都是常態。
\end{restatable}
% 附錄：同一個編號、同一段文字
\cw*
\begin{proofof}{thm:cw}
看特徵函數。
\end{proofof}
\end{cs}

\begin{cs}
\begin{lemma}[number=7.4] 抄書的編號，
之後的自動編號不受影響。
\end{lemma}
\begin{remark*} 不編號的版本加星號。
\end{remark*}
\begin{axiom}[選擇公理] 非空集合的
乘積非空。
\end{axiom}
\end{cs}
\csnote{全部的種類：\opt{theorem lemma proposition corollary claim definition assumption example remark}，另有 \opt{axiom conjecture property conclusion remarks notationthm problem}；都有不編號的 \opt{*} 版本。}

\begin{cssrc}{note}
\akizumiNewTheorem{tool}{Tool}{工具}{道具}{도구}[prefix=T]  % 工具 T1
\akizumiSetTheorem{assumption}{counter=own, prefix=A}       % 假設 A1
\listoftheorems                                             % 定理列表
\end{cssrc}

\section{提示框}
\begin{cs}
\begin{keyidea}
級數收斂，就是尾巴越來越小。
\end{keyidea}
\begin{tip}[根值判別法]
看 $\sqrt[n]{a_n}$ 的極限。
\end{tip}
\begin{warning}
$a_n \to 0$ 不保證收斂。
\end{warning}
\end{cs}

\begin{cs}
\begin{proofidea}
先比較，再用 Cauchy 準則。
\end{proofidea}
\begin{note} 也有 pitfall。
\end{note}
\begin{setup}
設 $(a_n)$ 為實數數列。
\end{setup}
\begin{remark}[\optional 選讀]
\code{\optional} 也能放在章節標題前。
\end{remark}
\end{cs}

\begin{cs}
\begin{chapterguide}
\guidehead{在談什麼} 級數的收斂。
\guidehead{先備知識} 數列極限。
\end{chapterguide}
\end{cs}

\begin{cssrc}{note}
\akizumiNewCallout{keypoint}{Key point}{學習要點}{}{}[numbered=chapter]
\begin{keypoint}\label{kp:a} ... \end{keypoint}  % 學習要點 3.1，可以 \cref
\end{cssrc}

\section{習題與解答（講義、書）}
\begin{cs}
\begin{exercise}[聯集]\label{ex:u}
求 $[-2,1) \cup (0,3]$。
\end{exercise}
\begin{solution}
$[-2,3]$。
\end{solution}
\begin{example-problem}[title=Two Sum,
  source=LeetCode 1]
找兩個數，和為 target。
\begin{parts}
\item 暴力法 \item 雜湊表
\end{parts}
\end{example-problem}
\end{cs}

\begin{cs}
\begin{exercises}[title]
\item 證明 $\sqrt{2}$ 是無理數。
\item 計算 $\sum_{k=1}^n k$。
  \begin{solution}
  $\frac{n(n+1)}{2}$
  \end{solution}
\end{exercises}
\end{cs}

\begin{cssrc}{book}
\documentclass[solutions=end]{akizumi-book}  % 解答收到書末
\printsolutions                              % 或印在這裡（例如每章最後）
\end{cssrc}

\begin{cssrc}{hw}
\documentclass[lang=zh, kind=report, header=name-id]{akizumi-hw}
\akizumisetup{course=常微分方程, number=3, due={10 月 30 日}}
\begin{document}
\maketitle
\problemset{B}                        % 之後的題號是 B1、B2…
\begin{problem}[一階線性方程]
  \begin{parts}  \item 求通解。  \item 解 $y(0)=1$。  \end{parts}
  \begin{solution}  積分因子為 $e^{2x}$。  \end{solution}
\end{problem}
\begin{problem}[label=B2$'$, title=緊緻性]
  \begin{parts} \item \begin{subparts} \item ... \end{subparts} \end{parts}
\end{problem}
\inputproblem[柯西數列]{problems/B3}  % 一題一個檔
\end{cssrc}
\csnote{\opt{kind=}：\opt{homework assignment exercise report}（標題的字）；\opt{header=name-id}：頁首印「作者 · 學號」；\opt{problembreak}：每題換頁；\opt{solutions=false} 隱藏解答。}

\section{數學}\nobreak
\setlength{\tabcolsep}{4pt}
\begin{tabular*}{\linewidth}{@{\extracolsep{\fill}}l l l l@{}}
\toprule
\mrow{\vb{x}\T\vb{A}\vb{x}} & \mrow{\va + \vb{\mu}} \\
\mrow{\norm{x}_2,\ \abs{x}} & \mrow{\inner{x, y}} \\
\mrow{\norm*{\frac{x}{2}}} & \mrow{\norm[\big]{x}} \\
\mrow{\set{x \given x > 0}} & \mrow{\floor{x},\ \ceil{x}} \\
\mrow{\E[X],\ \Var[X]} & \mrow{\Prob{A \given B}} \\
\mrow{\R,\ \N,\ \Z,\ \Q,\ \C} & \mrow{\mathbb{1},\ \mathscr{F}} \\
\mrow{\argmin_{\theta} L} & \mrow{\int f \dd x} \\
\mrow{\Normal{\mu}{\sigma^2}} & \mrow{X \indep Y,\ X_i \iid F} \\
\mrow{\tr,\ \rank,\ \spn,\ \im} & \mrow{\cvec[p]{a, b}} \\
\bottomrule
\end{tabular*}
\csnote{\opt{\textbackslash va}…\opt{\textbackslash vZ} 要 \opt{letters=bold}；\opt{delims=auto} 讓括號不加星號就自動變大；\opt{units} 載入 siunitx。}

\section{參考文獻、圖}
\begin{cssrc}{note}
\documentclass[bib=refs.bib, bibstyle=authoryear]{akizumi-note}
\citet{lin2026} 提出……；後續研究 \citep{lin2026, sutton2018}。
\printbibliography
\begin{figure}[t]
  \centering
  \begin{subfigure}{0.45\linewidth}
    \includegraphics[width=\linewidth]{figures/a} \caption{訓練}
  \end{subfigure}
  \caption{損失曲線}\label{fig:loss}
\end{figure}
\end{cssrc}

\section{程式碼}
\begin{cs}[0.47]
\begin{codeblock}[lang=Python, title=f.py]
def f(z, k=64):
    u = randn(z.shape[1], k)
    return (z @ u).mean()
\end{codeblock}
行內：\code{loss.backward()}、
\code[cpp]{std::vector<int> v;}
\end{cs}

\begin{cs}[0.47]
\begin{codeblock}[numbered, label=lst:fit,
  lang=python, numbers=false]
model.fit(x, y)
\end{codeblock}
見\cref{lst:fit}。
\begin{terminal}[title=zsh]
$ python main.py
3
\end{terminal}
\iosample{1 2 \\ 3 4}{3 \\ 7}
\end{cs}

\begin{cssrc}{note}
\codedir{src/}                               % \inputcode 也在 src/ 找檔案
\inputcode[lang=cpp, lines=10-24]{main.cpp}  % 讀檔，只印第 10-24 行
\inputcode[output, title=輸出]{run.txt}      % 程式輸出：白底、無行號
\listofcodes                                 % 所有編號過的程式碼
\end{cssrc}
\csnote{\opt{code=minted} 改用 Pygments 上色，寫法完全一樣；\opt{verbatim=styled} 讓 \opt{verbatim} 也有灰底。}

\section{演算法、符號表}
\begin{cs}
\begin{algorithm}[H]
\caption{SIGReg 訓練}\label{alg:s}
\begin{algorithmic}[1]
\Require 批次 $\{\vb{o}_i\}$，權重 $\lambda$
\State $\vb{Z} \gets f_\theta(\vb{o})$
\State $\theta \gets \theta - \eta
  \nabla_\theta \mathcal{L}$
\end{algorithmic}
\end{algorithm}
\end{cs}

\begin{cs}
\begin{notation}
\notn{$\vb{z}_t$}{時間 $t$ 的潛在狀態}
\notn{$\sg(\cdot)$}{stop-gradient}
\end{notation}
\end{cs}

\section{TikZ}
\begin{cs}
\begin{tikzpicture}
\node[akizumi/box] (a) {encoder};
\node[akizumi/box, right=1cm of a]
  (b) {$\vb{z}$};
\draw[akizumi/arrow] (a) --
  node[above, akizumi/note] {sg} (b);
\end{tikzpicture}
\end{cs}

\begin{cs}
\begin{tikzpicture}[x=0.9cm]
\foreach \v [count=\i from 0] in
    {3, 1, 4, 1}
  \node[akizumi/cell] at (\i, 0) {\v};
\node[akizumi/hlcell] at (2, 0) {4};
\node[akizumi/index] at (2, -0.7) {i};
\end{tikzpicture}
\end{cs}

\begin{cs}
\begin{tikzpicture}[akizumi/automaton]
\node[state, initial] (q0) {$q_0$};
\node[state, accepting, right=of q0]
  (q1) {$q_1$};
\path[->] (q0) edge["$a$"] (q1)
  (q1) edge[loop above] node {$b$} ();
\end{tikzpicture}
\end{cs}
\csnote{流程圖 \opt{akizumi/start} \opt{process} \opt{decision} \opt{io} \opt{group}；賽局樹 \opt{akizumi/game/...}；電流 \opt{\textbackslash pic\{akizumi/current in\}}。}

\section{邏輯（\texttt{\textbackslash usepackage\{akizumi-logic\}}）}
\begin{cs}
\begin{fitch}
\hypo{1}{A \limp B} \by{\RulePremise}{}
\open
\hypo{2}{A} \by{\RuleA[\RuleImpI]}{}
\have{3}{B} \by{\RuleImpE}{1,2}
\close
\have{4}{A \limp B} \by{\RuleImpI}{2-3}
\end{fitch}
\end{cs}

\begin{cs}
\begin{truthtable}{cc|ccc}
A & B & A & \main{\land} & B \\ \midrule
\tvrow{TT|T!TT}
\tvrow{TF|T!FF}
\end{truthtable}
\begin{calc}
\fn{sum}\;[1,2]
\step{=}{def.\ of $\fn{sum}$}
1 + \fn{sum}\;[2]
\end{calc}
\end{cs}
\csnote{也有 ebproof 的 \opt{prooftree}、\opt{\textbackslash rulename}、\opt{funcdef}、\opt{\textbackslash sem}；\opt{connectives=classic} 把 $\limp$ 換成 $\supset$。}

\section{其他小工具}
\begin{cs}
\todo{補推導}

\muted{灰色的次要文字}

課名\akizumisep 第 3 講

顏色：\textcolor{accent}{accent}、
\textcolor{warn}{warn}、\textcolor{mute}{mute}
\end{cs}

\begin{cssrc}{note}
\sidenote{邊欄註記}            % layout=margin 才在邊欄，否則是註腳
\blindfootnote{本講義改編自…}  % 沒有編號的註腳
\index{分割!網目}              % index=true 才有索引
\nota{$\norm{x}$}{歐氏範數}    % 符號列表
\printnotation  \printindex
\chapterdir{basic-op}          % 之後的圖、\input、程式碼先找 basic-op/
\end{cssrc}

\begin{cssrc}{中文報告}
\documentclass[lang=zh, par=indent, secnum=zh, 13pt, kai,
  text=times]{akizumi-note}
\akizumisetup{subtitle={食品安全風險溝通之案例研究},
  info={學號=B10303083, 系級=資工四}}     % 印在標題下方
\end{cssrc}

\section{簡報、考卷、書、論文（只列原始碼）}
\begin{cssrc}{slides}
\documentclass[aspectratio=169]{beamer}
\usetheme[progress]{Akizumi}        % 底部進度條；也接受 lang=zh、tikz=full
\begin{document}
\begin{frame}[plain,noframenumbering]\titlepage\end{frame}
\section{動機}                      % 每節前自動插入「Part N」
\begin{frame}[fragile]{程式碼要 fragile}
  \q{這是討論問題嗎？}              % 藍色的 Q:
\end{frame}
\standoutframe{Questions?}          % 整頁藍底一句話
\appendix  \dividerframe{Appendix}  % 之後的頁不算進總頁數
\end{cssrc}

\begin{cssrc}{exam}
\documentclass[lang=zh]{akizumi-exam}       % solutions=true 印出解答
\begin{document}
\begin{center}\gradetable[h][questions]\end{center}
\begin{questions}
\question[10] 求 $y'' + y = 0$ 的通解。
  \begin{solution} $y = c_1\cos x + c_2\sin x$。 \end{solution}
\question \begin{parts} \part[5] ...  \part[10] ... \end{parts}
\end{questions}
\end{cssrc}

\begin{cssrc}{book}
\documentclass[lang=zh, openright, chaptername={第 * 講}]{akizumi-book}
\frontmatter  \maketitle  \preface ...  \tableofcontents
\mainmatter
\part[sub=Foundations]{基礎}
\chapter[sub=Basic Operations, short=運算]{基本運算}
\chaptertoc                               % 本章小目錄
\backmatter  \printbibliography
\end{cssrc}

\begin{cssrc}{paper}
\usepackage{akizumi-paper}                 % 放在序言最後
\newcommenter{jun}{Jun}{teal}
\jun{這個界夠緊嗎？}  \jun*{改寫過的句子}  % 共同作者的註解、修改
如 \eg 表中 \best{0.91}、\second{0.88}；\anon{我們先前的工作}
\end{cssrc}

\end{document}
